\documentclass[border=10pt]{standalone}
\usepackage{fontspec}
\setmainfont{Noto Serif}[Ligatures=TeXOff]
\XeTeXgenerateactualtext=1
\newfontfamily\bengalifont{Noto Serif Bengali}[Script=Bengali, NFSSFamily=bnfam]
\usepackage[Latin,Bengali]{ucharclasses}
% group-free font switching (as in header-common.tex): robust inside braces, cells and quotes
\makeatletter
\def\BN@saved{\rmdefault}
\def\BN@bn{bnfam}
\newcommand\BNin{\edef\BN@tmp{\f@family}\ifx\BN@tmp\BN@bn\else\let\BN@saved\BN@tmp\fi\fontfamily{bnfam}\selectfont}
\newcommand\BNout{\fontfamily{\BN@saved}\selectfont}
\makeatother
\setTransitionsFor{Bengali}{\BNin}{\BNout}
% Typographic quotes and dashes are not in a ucharclasses class (class 0), so a Bangla word right
% after “ or ‘ would not trigger the Bangla font. Treat them as Latin so that every switch fires.
\XeTeXcharclass"2018=\XeTeXcharclass`A
\XeTeXcharclass"2019=\XeTeXcharclass`A
\XeTeXcharclass"201C=\XeTeXcharclass`A
\XeTeXcharclass"201D=\XeTeXcharclass`A
\XeTeXcharclass"2013=\XeTeXcharclass`A
\XeTeXcharclass"2014=\XeTeXcharclass`A
\XeTeXcharclass"2026=\XeTeXcharclass`A
\usepackage{newunicodechar}
\newunicodechar{।}{{\bengalifont\char"0964}}
\newunicodechar{॥}{{\bengalifont\char"0965}}

\usepackage[table]{xcolor}
\usepackage{booktabs,tabularx,array}
\definecolor{bapilrule}{HTML}{1F2937}
\definecolor{bapilstripe}{HTML}{F3F4F6}
\renewcommand{\arraystretch}{1.28}
\setlength{\tabcolsep}{5pt}

\begin{document}
\begin{minipage}{22cm}
{\large\bfseries সারণি ১৫।}\ {\large দায়িত্ব-ম্যাট্রিক্স (R দায়িত্বপ্রাপ্ত, A জবাবদিহিকারী, C পরামর্শপ্রাপ্ত, I অবহিত)}\par\vspace{7pt}
\small
\rowcolors{2}{bapilstripe}{white}
\begin{tabularx}{\linewidth}{@{}>{\hsize=2.658\hsize\raggedright\arraybackslash}X>{\hsize=0.709\hsize\raggedright\arraybackslash}X>{\hsize=0.709\hsize\raggedright\arraybackslash}X>{\hsize=0.709\hsize\raggedright\arraybackslash}X>{\hsize=0.709\hsize\raggedright\arraybackslash}X>{\hsize=0.709\hsize\raggedright\arraybackslash}X>{\hsize=0.797\hsize\raggedright\arraybackslash}X@{}}
\arrayrulecolor{bapilrule}\toprule
\rowcolor{white}\textbf{সিদ্ধান্ত বা কার্যক্রম} & \textbf{স্টিয়ারিং বোর্ড} & \textbf{ERRC} & \textbf{TAB} & \textbf{পরিচালক} & \textbf{ইউনিট} & \textbf{স্বাগতিক সংস্থা} \\
\midrule
বার্ষিক পরিকল্পনা ও বাজেট & A & C & C & R & C & I \\
সনদ সংশোধন & A & C & I & R & I & I \\
G0-তে স্তর যাচাই & I & I & C & A & R & C \\
G1-এ নকশা অনুমোদন & I & C & A & R & R & C \\
G2-তে পাইলট অনুমোদন (স্তর T2) & C & A & C & R & R & C \\
G4-এ গ্রহণের সিদ্ধান্ত & A & C & C & R & C & A \\
গুরুতর ঘটনার পর স্থগিতকরণ & I & A & C & R & R & C \\
মডেল ও নির্ভরতা নির্বাচন & I & I & A & C & R & I \\
ডেটা-বিনিময় চুক্তি & I & C & C & A & R & C \\
স্বচ্ছতা নিবন্ধন প্রকাশ & I & C & I & A & R & I \\
জ্যেষ্ঠ কর্মী নিয়োগ & A & I & I & R & I & I \\
ফলাফলের জনসমক্ষে প্রচার & I & C & I & A & R & C \\
\bottomrule
\end{tabularx}
\par\vspace{5pt}{\footnotesize\itshape G4-এ স্টিয়ারিং বোর্ড ও স্বাগতিক সংস্থা যৌথভাবে জবাবদিহিকারী: পরিচালনাগত সিদ্ধান্তের মালিক সংস্থা। A দিয়ে BAPIL-এর অভ্যন্তরীণ জবাবদিহি বোঝানো হয়েছে; BAPIL-এর কোনো আইনগত সত্তা নেই বলে চুক্তি ও ডেটা-বিনিময় চুক্তি স্বাগতিক প্রতিষ্ঠান স্বাক্ষর করে।}
\par\vspace{3pt}{\scriptsize\color{gray} উৎস: BAPIL টেকনিক্যাল ব্লুপ্রিন্ট, সংস্করণ ১.০ (২০২৬)। CC BY 4.0।}
\end{minipage}
\end{document}
